\documentclass[12pt]{article}
\usepackage[utf8x]{inputenc}
\usepackage[english, russian]{babel}
\usepackage{multicol}
\usepackage{amsthm}
\usepackage{unicode-math}
\usepackage[%
    left=1in,%
    right=1in,%
    top=1in,%
    bottom=1in
]{geometry}%
\usepackage{graphicx}
\graphicspath{{img/}}

\begin{document}
\begin{multicols}{2}
\noindentв него входила только одна тригонометрическая функция угла $\alpha$. Для этого представим коэффициент трения k как тангенс некоторого угла $\phi$:
\begin{center}
    $\tan\phi=k$; $\phi=\arctan k$
\end{center}
\newline
\begin{center}
    $\sin\phi=\frac{k}{\sqrt{1+k^2}}$; $\cos\phi=\frac{1}{\sqrt{1+k^2}}$
\end{center}


Так как $k$ задано, угол ф можно считать известным. Подставляя в (4) $\sin\phi$ вместо $k$ и умножая числитель и знаменатель на $\cos\alpha$, получим
\newline
$F=\frac{P\tan\phi}{\tan\phi\sin\alpha+\cos\alpha}=\frac{P\sin\phi}{\sin\phi\sin\alpha+\cos\phi\cos\alpha}=\frac{P\sin\alpha}{\cos(\alpha-\phi)}=\frac{Pk}{\sqrt{1+k^2}\cos(\alpha-\phi)}$
Теперь ясно, что сила F будет минимальной тогда, когда $\cos(\alpha - \phi) = 1$, тo есть
$$\alpha_{min} = \phi = \arctan k$$
При этом минимальное значение
силы $F$ равно $$F_{min}=\frac{kP}{\sqrt{1+k^2}}$$


Из решения этой задачи можно сделать практический вывод: когда необходимо везти на санях груз по дороге с большим коэффициентом трения (например, если дорога посыпана песком), нужно тянуть сани за короткую веревку. Если же коэффициент трения мал, веревка должна быть длинной. Объясните это. Задача 3. Точки 1 и 2 движутся по осям $x$ и $y$ к началу координат (рис. 2). В момент $t_0 = 0$ точка 1 находится на расстоянии $s_1 = 10$ см,

\includegraphics[scale=0.7]{img2.png}
\newline
\caption{Рис. 2}
\columnbreak
\newline
а точка 2 - на расстоянии $s_q = 5$ см от начала координат. Первая точка движется со скоростью $v_1 = 2$ см/с, а вторая - со скоростью $v_2 = 4$ см/с. Каково наименьшее расстояние между ними? В момент $t$ расстояние между точками равно
$$d=\sqrt{(s_1-v_1t)^2+(s_2-v_2t)^2}=$$
$$=\sqrt{(10-2t)^2+(5-4t)^2}=\sqrt{20t^2-80t+125}$$


Под корнем стоит квадратный трехчлен, из которого можно легко выделить полный квадрат. Действительно,
$$20t^2-80t+125=5[(2t-4)^2+9]$$


Ясно, что минимум этого выражения, а значит, и расстояния $d$ будет при $t_2$, откуда
$d_{min}=\sqrt{45}$

\includegraphics[scale=0.7]{img3.png}
\newline
\caption{Рис. 3}
\end{multicols}
\end{document}